\section{The rules dialect in full}
\label{app:dialect}

This appendix states the dialect of \S\ref{sec:rules} at the length required to
re-implement it. Table~\ref{tab:dialect} is the index: it lists every
consequential choice against the engine switch that changes it, and the
subsections that follow give each choice in full. The two files that carry the
definitions are \filepath{engine/src/fish.hpp} (deck, indexing, the
\texttt{Rules} record, the legality predicate) and \filepath{engine/src/game.hpp}
(the driver: declaration rounds, turn handling, the forced endgame, the action
cap). Where the text below names a field it is a member of \texttt{Rules}; where
it names a routine it is a member of the driver.

\begin{table}[ht]
\centering
\caption{The rules dialect used for every experiment in this paper unless
stated otherwise, and the engine switch that changes each choice. The switches
are command-line flags of the \texttt{fish} driver; the underlying fields are
members of \texttt{Rules} in \texttt{engine/src/fish.hpp}. Rows marked ``modelling
choice'' are not fixed by any published statement of the rules.}
\label{tab:dialect}
\small
\begin{tabular}{@{}p{0.29\linewidth}p{0.46\linewidth}p{0.19\linewidth}@{}}
\toprule
Choice & Value used & Switch \\
\midrule
Deck composition & 54 cards in nine half-suits (low and high band of each suit,
plus the four eights and two jokers); nine cards per player &
\texttt{-{}-sets=8} gives the 48-card, eight-half-suit form \\
Ask legality & Live half-suit; asker holds another card of it and not the named
card; target is an opponent and still holds cards & --- \\
Turn transfer & A successful ask retains the turn; a miss passes it to the
target & --- \\
Out-of-turn declarations & Permitted at any moment, including during an
opponent's turn & \texttt{outOfTurnDeclare} (\texttt{-{}-no-out-of-turn}) \\
Cardless players may declare & Permitted &
\texttt{cardlessMayDeclare} (\texttt{-{}-no-cardless-declare}) \\
Incorrect declaration & Awards the half-suit to the opposing team & --- \\
Declaration aftermath & The six cards leave play and the turn returns to the
player who held it & --- \\
Cardless turn-holder & Chooses which live teammate receives the turn &
--- \\
Forced endgame (whole team cardless) & The live team declares every remaining
half-suit, sharing only willingness, over an eight-level ladder
$\{0.995, 0.98, 0.95, 0.90, 0.80, 0.65, 0.50, \text{best guess}\}$ &
\texttt{-{}-legacy} gives the two-level ladder $\{0.38, \text{best guess}\}$ \\
Simultaneous declarations & Lowest seat index declares first (a modelling
choice; alternatives measured in \S\ref{app:dialects}) &
\texttt{declArbitration} (\texttt{-{}-arb=low\textbar high\textbar turn}) \\
Action cap & 400 asks (360 under the v0.3 dialect), any residue adjudicated
neutrally by majority physical holding, ties to the holder of the lowest card &
\texttt{-{}-maxasks} \\
Conventions & No prearranged signalling conventions; team communication is
limited to the willingness bit described above & --- \\
\bottomrule
\end{tabular}
\end{table}

\subsection{Players, teams, and the deck}

Six players occupy seats $0,\dots,5$. Teams alternate: seats $\{0,2,4\}$ form
team~0 and seats $\{1,3,5\}$ form team~1, so every player's two neighbours are
opponents. The deck is 54 cards partitioned into $N_H = 9$ half-suits of six
cards. Half-suits are indexed $h \in \{0,\dots,8\}$ and positions within a
half-suit $i \in \{0,\dots,5\}$; a card is the integer
\[
  c \;=\; 6h + i, \qquad h = \lfloor c/6 \rfloor, \qquad i = c \bmod 6 .
\]
The ordering is fixed and matches the v0.3 engine card for card, so that ported
policies and cross-engine checks line up:

\begin{center}
\small
\begin{tabular}{@{}clp{0.46\linewidth}@{}}
\toprule
$h$ & Half-suit & Cards at positions $i = 0,\dots,5$ \\
\midrule
0 & Low Spades    & 2S, 3S, 4S, 5S, 6S, 7S \\
1 & High Spades   & 9S, 10S, JS, QS, KS, AS \\
2 & Low Hearts    & 2H, 3H, 4H, 5H, 6H, 7H \\
3 & High Hearts   & 9H, 10H, JH, QH, KH, AH \\
4 & Low Diamonds  & 2D, 3D, 4D, 5D, 6D, 7D \\
5 & High Diamonds & 9D, 10D, JD, QD, KD, AD \\
6 & Low Clubs     & 2C, 3C, 4C, 5C, 6C, 7C \\
7 & High Clubs    & 9C, 10C, JC, QC, KC, AC \\
8 & Eights and Jokers & 8S, 8H, 8D, 8C, red joker, black joker \\
\bottomrule
\end{tabular}
\end{center}

Even $h$ are low bands, odd $h$ are high bands, and $\lfloor h/2 \rfloor$ is the
suit in the order spades, hearts, diamonds, clubs; $h = 8$ is the mixed
half-suit. Hands are 64-bit masks with bit $c$ set when the seat holds card $c$,
so a half-suit is the mask $\mathtt{0x3F} \ll 6h$. The code identifier for a
half-suit is \texttt{set}; the term is used nowhere else in this paper.

With \texttt{deckSets} $=8$ (the \texttt{-{}-sets=8} switch) the ninth half-suit
is removed, giving the 48-card, eight-half-suit form with eight cards per
player; all other rules are unchanged.

\subsection{The deal and the opening turn}

The deck is shuffled by a Fisher--Yates pass driven by a splitmix64 generator
seeded from the deal seed, and dealt in contiguous blocks of $54/6 = 9$ cards
(eight cards under \texttt{deckSets} $=8$) to seats $0,\dots,5$ in order. The
initial deal is retained separately as ground truth and is never exposed to a
policy. A dealer seat is then drawn uniformly and the opening turn is given to
the seat clockwise of the dealer. Because the shuffle is a pure function of the
seed, a deal is reproducible from its seed alone, which is what makes the
duplicate-block protocol of \S\ref{sec:protocol} possible.

\subsection{Asking}
\label{app:dialect-ask}

Let $\mathcal{H}$ be the set of live half-suits, $\mathrm{team}(p)$ the team of
seat $p$, and $n_p$ the public card count of seat $p$. Seat $a$ asking seat $t$
for card $c$ is legal exactly when all five of the following hold:

\begin{enumerate}
  \item $h(c) \in \mathcal{H}$: the half-suit has not yet been claimed;
  \item $\mathrm{team}(t) \neq \mathrm{team}(a)$: the target is an opponent;
  \item $n_t > 0$: the target still holds cards;
  \item $a$ does not hold $c$;
  \item $a$ holds at least one card of $h(c)$ --- necessarily a card other than
        $c$, by the previous condition.
\end{enumerate}

Two properties of this predicate matter for the rest of the paper. First, the
only private input is the asker's own hand: conditions 1--3 are functions of the
public state alone, and conditions 4--5 are functions of the public state and
the asker's hand. An observer watching a legal ask therefore learns a certificate
about the asker's hand without needing to see it, which is the C5 constraint of
\S\ref{sec:belief}. Second, the predicate is evaluated identically on the
ground-truth hand inside the driver and on the observer's reconstruction, so the
two cannot disagree.

The target must surrender the card if they hold it, and no other transfer is
possible. On a hit the asker keeps the turn; on a miss the turn passes to the
target. The ask, the identity of asker and target, the named card, the outcome,
and the resulting card counts are all appended to a public event history.

If the seat to move has no legal ask at all --- which happens only when every
card it holds lies in a half-suit whose other five cards it also holds --- it
must declare instead, and the driver requires a declaration of one such
half-suit from that seat.

\subsection{Declarations}
\label{app:dialect-declare}

A declaration names a live half-suit $h$ and an allocation
$A: \{0,\dots,5\} \to \{0,\dots,5\}$ assigning each position $i$ of $h$ to the
seat claimed to hold card $6h+i$. Three conditions are checked before the
declaration is applied: the half-suit must still be live; every named seat must
belong to the declarer's own team; and, when \texttt{outOfTurnDeclare} is
disabled or \texttt{cardlessMayDeclare} is disabled, the declarer must
respectively hold the turn or hold at least one card. Under the dialect studied
both switches are enabled, so any seat may declare at any moment whether or not
it holds cards and whether or not it holds cards of the half-suit concerned.

Resolution is all-or-nothing. The declaration is correct when, for every
$i$, seat $A(i)$ actually holds card $6h+i$; a single misassignment --- including
one that names the right team but the wrong teammate --- makes the whole
declaration incorrect. A correct declaration awards $h$ to the declarer's team;
an incorrect one awards $h$ to the opposing team. In both cases the six cards
are removed from every hand, $h$ leaves $\mathcal{H}$, the public card counts are
updated, and the turn returns to the seat that held it before the declaration.
The declarer's own stated confidence is recorded in the event history as a
diagnostic and is never used to resolve anything.

The driver polls for declarations before every ask, in a loop that repeats while
any seat still offers a declaration, so several half-suits can be claimed between
two consecutive asks.

\subsection{Arbitration between simultaneous declarations}
\label{app:dialect-arb}

Because any seat may declare at any moment, the poll can find more than one
willing declarer, and the rules as published do not say who acts first.
\texttt{declArbitration} selects the scan order over seats and the first willing
seat in that order declares:

\begin{itemize}
  \item \texttt{0} (default, \texttt{-{}-arb=low}): ascending seat index
        $0,1,\dots,5$;
  \item \texttt{1} (\texttt{-{}-arb=high}): descending seat index
        $5,4,\dots,0$;
  \item \texttt{2} (\texttt{-{}-arb=turn}): ascending from the current
        turn-holder, $g, g+1, \dots$ modulo six.
\end{itemize}

Selecting the most confident proposer is deliberately not offered. Confidence is
private to a seat, and comparing confidences across seats would move information
between opponents as well as teammates; it is the same leak that the forced
endgame was corrected to remove (\S\ref{app:dialect-vthree}). E17 reports the primary
head-to-head result under all three orders (\S\ref{app:dialects}).

\subsection{Cardless players and the turn}
\label{app:dialect-cardless}

A seat with no cards leaves the asking cycle: it cannot ask, and condition~3 of
\S\ref{app:dialect-ask} makes it an illegal target. It remains part of the game
in two ways --- it can be named in a teammate's declaration as the holder of a
card, which is impossible while it is genuinely empty and therefore constrains
the allocations a teammate may name, and, with \texttt{cardlessMayDeclare}
enabled, it may itself declare.

The turn can reach a cardless seat only immediately after a declaration removed
its last cards. That seat then chooses which of its live teammates receives the
turn; the choice is a policy decision, and an invalid choice is replaced by the
lowest-indexed live teammate. If the seat has no live teammate, its whole team is
cardless and the forced endgame begins.

\subsection{The forced endgame}
\label{app:dialect-endgame}

When one team holds no cards at all, no legal ask remains for either side, and
the live team must declare every remaining half-suit. Teammates may negotiate
openly but may exchange nothing except willingness to declare, so the driver
implements the selection as a threshold sweep rather than a comparison of
confidences. The willingness ladder is
\[
  \Theta \;=\; \{\,0.995,\; 0.98,\; 0.95,\; 0.90,\; 0.80,\; 0.65,\; 0.50,\;
  \text{best guess}\,\},
\]
eight levels held in \texttt{forcedTh}, the last entry (stored as $-1$) meaning
that somebody must declare and should submit their best guess. For each
threshold $\theta \in \Theta$ in descending order, the driver scans the live
team's seats over the remaining half-suits and asks each seat whether it is
willing to declare that half-suit at confidence at least $\theta$; the first
seat that says yes declares, the ladder restarts at its top, and the sweep
repeats until every half-suit is resolved. Only the willingness bit crosses
between seats. A consequence of restarting the ladder is that early, safer
declarations resolve cards and thereby sharpen the allocations available for the
later ones.

Declarations in the forced endgame resolve exactly as in
\S\ref{app:dialect-declare}, including the possibility of misdeclaring and
handing a half-suit to the cardless team. They are recorded separately from
voluntary declarations in the event history. If every seat refuses at every
level --- reachable only with a policy that declines its own best guess --- the
remaining half-suits are awarded to the team physically holding the majority of
their cards.

Under \texttt{-{}-legacy} the ladder is the two-level
$\{0.38,\ \text{best guess}\}$ of the v0.3 dialect.

\subsection{Two further differences from the v0.3 dialect}
\label{app:dialect-vthree}

\S\ref{sec:rules-v03} lists the two v0.3 differences that bear on the reported
results --- turn-only declarations and the 360-ask cap. The other two concern
the two teammate-to-teammate channels of \S\ref{app:dialect-conventions}, and
are recorded here.

\paragraph{The successor of a cardless turn-holder was fixed.} In the v0.3
engine \cite{fishbot03} the turn passed automatically to the live teammate
holding the most cards. In the dialect studied the cardless seat chooses
(\S\ref{app:dialect-cardless}). The v0.3 rule is a function of the public card
counts alone and so conveys nothing between teammates; the choice made here is a
policy decision and is therefore the second of the two channels listed in
\S\ref{app:dialect-conventions}. What it can carry is bounded by its range: one
selection among at most two live teammates, made only immediately after a
declaration emptied the chooser's hand.

\paragraph{The forced endgame selected the declarer by comparing private
confidences.} v0.3 gave the declaration to whichever live seat stated the
highest confidence in a remaining half-suit. Ordering seats by confidence
conveys more than willingness, because the ordering is itself a statistic of the
hands being compared. The threshold sweep of \S\ref{app:dialect-endgame}
replaces that comparison with a descending sequence of thresholds shared in
advance, so that only a willingness bit crosses between seats. The same argument
is why confidence ranking is not offered as an arbitration order in
\S\ref{app:dialect-arb}.

\subsection{The action cap and adjudication}

The driver counts asks and stops the game once \texttt{maxAsks} is reached: 400
under the dialect studied, 360 under \texttt{-{}-legacy}, and settable with
\texttt{-{}-maxasks}. Awarding the unresolved residue to either side would bias
the outcome towards whichever policy stalls more readily, so with
\texttt{adjudicateOnLimit} enabled each remaining half-suit is awarded to the
team physically holding the majority of its six cards, ties going to the team
holding the lowest-indexed card of that half-suit. The game result records
whether the cap was reached, and that incidence is reported alongside every
experiment in this paper. It is \numLimitRate{} in E1, E3, E7, E10 and E13; the
exception is the deliberately unforced mirror configuration of E11
(\S\ref{sec:termination}), which is measured precisely because it does reach the
cap.

\subsection{Termination of a game and scoring}

The game ends when no live half-suit remains: every half-suit has been claimed
by a declaration, resolved in the forced endgame, or adjudicated at the cap.
Each of the nine half-suits is awarded to exactly one team, so the final scores
sum to nine and the team taking at least five wins. Under \texttt{deckSets} $=8$
the scores sum to eight, a level four-all result is possible, and the driver
awards such a game to team~0; that case cannot arise with nine half-suits.

\subsection{Communication and conventions}
\label{app:dialect-conventions}

No prearranged signalling conventions are available to any policy in this paper.
The only channel between teammates beyond the public event history is the
willingness bit of \S\ref{app:dialect-endgame} and the choice of successor made
by a cardless turn-holder. Every other quantity a policy computes --- ownership
marginals, allocation probabilities, declaration confidences --- is private to
the seat that computed it, and the information-safety audit of
\S\ref{sec:engine} checks that no policy is ever handed a quantity derived from
another seat's hand.
